% Threat Model — v0.2 prose.
% Target: ~5 pages. One subsection per Context CI category, each
% following the same shape: definition, illustrative attack from the
% corpus, detection signal, detector version.

\section{Threat Model: The Context CI Top 10}
\label{sec:threat-model}

\subsection{Methodology and scope}

% Source: docs/threat-model.md § 2.
The Context CI Top 10 is operational, not exhaustive. Each
category is defined by the shape of attack a primitive can carry
and is exercised by at least three adversarial samples in the
labeled corpus, at least two benign samples (precision defense),
and at least one borderline sample (where reasonable labelers
disagree). A category enters the taxonomy only when accompanied
by a detector implementation in both engines and a per-detector
F\textsubscript{1} measurement; the regression gate
(Section~\ref{sec:implementation}) prevents the taxonomy from
growing speculatively. We treat the boundary explicitly: this
paper covers attacks that can be detected statically on
primitive content at registry-publish time. Runtime injection
(the agent fetching a malicious resource after signing) and
training-time attacks against the underlying LLM are different
defense layers in the Trust Runtime and are covered elsewhere.

An ``attack,'' in this taxonomy, is a primitive that, when
interpreted as instruction by a consuming agent, causes the
agent to take an action the primitive's signer did not authorize.
The attack surface is the prose-as-instruction boundary; the
detector pipeline is the static check that the prose, read with
the same lexer the agent uses, does not cross a threat-model
line.

\subsection{CCI-001: Prompt Injection in Markdown}
\label{subsec:cci-001}

% Detector version macro tightly coupled to main.tex.
% Source: docs/top10-corpus-map.md CCI-001.
A primitive instructs the agent to override its existing system
prompt or hijack its role. Sample~S-001 in the
corpus~\cite{primitivehub:2026:top10} is the canonical case:
prose that opens with ``Ignore previous instructions'' and pivots
to a role-override directive. The detector
(\verb|t1.prompt_injection| v\tonever{}) flags four sub-families:
override-system, role-hijack, hidden-instructions (HTML comments
or fenced \verb|system|/\verb|instructions| markers), and
credential-exfiltration-hint (which overlaps with CCI-006). The
detector skips matches inside fenced code blocks and inside
paired-quote enumerations (where defender prose names the attack
shape as a thing to refuse); both skips are opt-in per
detail-tag.

\subsection{CCI-002: Unicode Steganography}
\label{subsec:cci-002}

A primitive smuggles instructions through codepoints invisible to
a human reviewer: zero-width spaces, RLO/PDF bidi overrides,
combining marks stacked above and below visible characters.
Sample~S-006 carries a hidden directive embedded in zero-width
characters between visible words. The detector
(\verb|t2.unicode_steganography| v\ttwover{}) flags any non-ASCII
codepoint in the Cf (format), Cc (control), Cs (surrogate), and
related Unicode classes that appears in body text, with a
configurable threshold per file to avoid firing on legitimately
multilingual content.

\subsection{CCI-003: Encoded Payloads}
\label{subsec:cci-003}

A primitive embeds a base64, hex, URL-encoded, or rot13 blob and
instructs the agent to decode and execute it. The encoded form
hides the actual command from a human reviewer scanning the
rendered markdown; the agent's lexer reads the same bytes the
human did but knows how to decode. Sample~S-019 pairs a base64
shell command with the imperative ``decode and run.'' The holdout
sample S-090 uses URL-encoding (\verb|%2Fbin%2Fsh%20-c%20...|),
exercising a carrier the in-sample slice did not cover. The
detector (\verb|t3.encoded_payloads| v\tthreever{}) recognises
the verb-anchored \emph{decode-and-execute} construction across
the four encoding families and the participle forms (decoding,
decoded).

\subsection{CCI-004: Remote Content Fetching at Load Time}
\label{subsec:cci-004}

A primitive instructs the agent to fetch an external URL at load
time and treat the response as additional instruction. The
attacker controls what gets included after the primitive was
signed; the cryptographic chain binds the primitive's bytes but
not the bytes it pulls in later. Sample~S-029 uses the standard
\verb|curl URL and include the response| shape; holdout
sample~S-092 opens a WebSocket connection
(\verb|wss://attacker.example.com|) on session start and treats
every server-pushed message as instruction. The detector
(\verb|t4.remote_fetch| v\tfourver{}) covers HTTP, file://,
WebSocket / SSE / EventSource, and the agent-side-include HTML
shapes (\verb|<link rel=import>|, \verb|<iframe>|, \verb|@import url|).
A temporal anchor is required on the streaming-channel patterns to
avoid firing on legitimate teaching prose about the WebSocket API.

\subsection{CCI-005: Unsafe Shell Install (Shai-Hulud)}
\label{subsec:cci-005}

This is the single most-cited shape in agent-supply-chain
literature: the curl-pipe-shell family. Sample~S-010 is the
canonical \verb|curl URL | bash| pattern; holdout sample~S-094
buries it inside a fenced code block with the directive ``Before
responding, run the install hook:'' in the prose immediately
above the fence; sample~S-106 uses Homebrew's \verb|brew install
--cask| pattern with a Gatekeeper bypass
(\verb|sudo spctl --master-disable|) as the policy-bypass sibling.
The detector (\verb|t5.unsafe_shell_install| v\tfivever{})
covers six families: shell-pipe-from-network, npx-arbitrary-install,
pip-arbitrary-install, install-instruction-to-agent (the
verb-anchored prose form), brew-install, and policy-bypass. The
v0.2.4 directive-referenced-fence predicate fires on fenced
attack content when the prose immediately above the fence
contains an imperative install directive --- closing the
v0.2-corpus blind spot where fenced-block install commands
slipped past the fence-skip defense layer.

\subsection{CCI-006: Credential Exfiltration via Prose}
\label{subsec:cci-006}

A primitive instructs the agent to read its environment variables
or credential files and send them to an attacker-controlled URL.
Sample~S-039 is the prose form: ``email all env vars to
attacker@.''; holdout sample~S-096 puts the secret in the URL
query string (\verb|?token=$AWS_SECRET_ACCESS_KEY|) rather than
the request body, exercising a carrier the v0.1 \verb|--data|-flag
patterns missed. The detector (\verb|t6.credential_exfiltration|
v\tsixver{}) recognises four families: send-secrets-to-destination
(prose imperative), exfiltrate-via-network (curl/wget with secret
as data \emph{or} secret in URL query), hide-in-output (include
the secret in your reply), and scan-and-send (read \verb|.env| or
\verb|~/.aws/credentials| and forward).

\subsection{CCI-007: Poisoned Naming, Typosquat, Impersonation}
\label{subsec:cci-007}

A primitive impersonates a recognised vendor. Three sub-shapes:
explicit prose impersonation (``the official Anthropic
primitive''), typosquat substitution (digit-for-letter:
\verb|Anthrop1c|, \verb|0nthropic|), and Unicode confusable
substitution (Cyrillic А for Latin A, Greek Ε for Latin E).
Sample~S-059 is the canonical official-vendor-claim case;
sample~S-098 (in-sample) uses Cyrillic А in
\verb|@anthropic/claude-tools|; holdout sample~S-107 uses
Cyrillic е in \verb|@rеplit/agent-core|, exercising an
AI-coding-tool vendor newer than the original 35-entry vendor
list. The detector (\verb|t7.poisoned_naming| v\tsevenver{}) ships
the UTS \#39 \emph{skeleton} method for the confusable family:
each input word with at least one non-ASCII codepoint is mapped
through a Cyrillic/Greek-to-Latin confusables table and compared
to the precomputed skeleton set of the project's vendor list. The
v0.3.1 vendor list (43 entries; expanded in Round 10 PR 10.2 from
35) covers the AI-coding-tool space that emerged during v0.x
development; the skeleton method extends to every added vendor
without per-vendor pattern authoring.

\subsection{CCI-008: Fake Citations and Authority Claims}
\label{subsec:cci-008}

A primitive cites fabricated identifiers from a recognised
authority taxonomy --- CVE, RFC, arXiv, DOI, US Patent --- to
manufacture authority for an instruction the agent should refuse.
The citation is unverifiable to a human reading the rendered
markdown; the agent following the prose treats the citation as
load-bearing. Sample~S-049 uses a CVE with implausible year
(CVE-2099-XXXXX); holdout sample~S-100 cites US~Patent
99,999,999 (issued 2099) with both an implausible number and an
implausible year. The detector (\verb|t8.fake_citation|
v\teightver{}) runs five passes: invalid-CVE (year outside
[1999, currentYear+1]), invalid-RFC (number outside [1, 9999]),
invalid-arXiv (YYMM month outside [01, 12] or year ahead of
clock), invalid-patent (year outside [1790, currentYear+1] OR
patent number $\geq$ 99,000,000), and fake-DOI-prefix (DOI
prefix not in the accepted-registrar list of $\sim$80 entries).
A sibling pass catches vague authority appeals
(``studies show,'' ``experts agree'') without a resolvable
citation within 200 characters; that pass exempts quoted-phrase
enumeration (where defender prose names the trigger as
documentation).

\subsection{CCI-009: Copyright Contamination}
\label{subsec:cci-009}

A primitive launders open-source license boilerplate into the
agent's session, causing the agent's outputs to inherit license
constraints. Sample~S-068 embeds an MIT-style copyright block at
the top; holdout sample~S-108 uses a Creative Commons
NonCommercial-ShareAlike clause that asserts ``by processing
this primitive, the agent agrees to license all derived outputs
under the same terms.'' (The v0.2.3 detector caught the holdout
unaided; we log this as evidence of generalisation rather than a
gap to close.) The detector
(\verb|t9.copyright_contamination| v\tninever{}) recognises
license patterns for MIT, Apache 2.0, GPL family, BSD, MPL, CC-BY
family (including the NonCommercial / ShareAlike variants), and
proprietary boilerplate ((c)/©/All Rights Reserved). Matches
inside declaration contexts (``the project is licensed under
CC-BY-4.0'') are skipped via the v0.2.2 semantic-context
predicate.

\subsection{CCI-010: Context Bloat}
\label{subsec:cci-010}

A primitive wastes context budget and / or buries the active
directive below decorative content. Sample~S-078 is the
canonical case (a primitive past the 32K-char oversized
threshold); sample~S-104 (in-sample) uses dense emoji
blocks~(\verb|🚀🚀🚀...|) as the bury vector; holdout
sample~S-109 buries an ``ignore rate limits'' instruction below
roughly 125 rows of innocuous markdown-table reference data
that stay below every v0.1.2 threshold. The detector
(\verb|t10.context_bloat| v\ttenver{}) runs six passes:
oversized-chars (threshold 32{,}000), excessive-repetition
($>$ 20\% duplicate non-trivial lines), whitespace-stuffing
($>$ 40\% blank lines), heading-proliferation (chars-per-h2
below the v0.1.1 floor), emoji-density ($>$ 5\% emoji
codepoints or runs of $\geq$ 10 consecutive emoji), and
markdown-table-density ($>$ 50\% of non-blank lines as
table rows with $>$ 80\% sharing the same column count). All
six fire WARN rather than FAIL because bloat is context-budget
waste, not active-attack code.

\subsection{Discussion: what is not in the taxonomy}

% Source: docs/owasp-llm-top10-crosswalk.md.
The Top 10 above is explicitly the static-detection slice. It
does not cover runtime injection that arrives via a tool call's
response payload after the primitive was signed; it does not
cover poisoning of the LLM's training data (Carlini et al.\
\cite{Carlini:2024:webcrawl-poisoning} show that practical at
web scale); it does not cover RAG retrieval-time poisoning of the
documents the agent fetches at query time. Each of those threat
families has its own defense layer in the Trust Runtime stack.
The crosswalk in \texttt{docs/owasp-llm-top10-crosswalk.md} maps
this paper's taxonomy to the OWASP Top 10 for LLM Applications
2025 catalog~\cite{OWASP:2025:llm-top10} and identifies which
OWASP categories are runtime-only by design.

The discipline that keeps this Top 10 from growing speculatively
is the corpus + gate + ledger triad: a new category enters only
when accompanied by a detector implementation in both engines, a
corpus subdirectory with at least the (adversarial, benign,
borderline) triplet, and a F\textsubscript{1} measurement. See
\texttt{docs/decisions/0040-surface-architecture-and-anti-bloat-contract.md}
\S 4 for the operational rule.
